Las simulaciones y su revisión permitieron establecer las siguientes conclusiones dentro del alcance de un monopolo inmerso en medios homogéneos a 3 GHz.
\begin{enumerate}
\item Se implementó una simulación en Python y Meep que permitió estimar la atenuación aparente en concreto, agua salada y tejido muscular. A 5 px/mm se obtuvieron 0.0380, 0.7012 y 0.4669 dB/mm, respectivamente, con los IC del Cuadro~\ref{tab:fdtd-base}. Las discrepancias teóricas fueron 1.31\,\%, 4.29\,\% y 5.12\,\%. El ajuste en vacío incumplió el criterio de calidad, lo que restringió su interpretación.
\item Se propuso un procedimiento reproducible de medición numérica basado en ventanas radiales declaradas, intervalos por bloques, comparación teórica, refinamiento y balance sobre un volumen identificado. La trazabilidad se conservó mediante archivos originales, códigos, identificadores y hashes. La incertidumbre del ajuste quedó diferenciada de la sensibilidad a la malla y del residuo energético.
\item El refinamiento de agua salada redujo la discrepancia de atenuación respecto de la teoría hasta aproximadamente 0.81\,\% a 40 px/mm. Las diferencias entre las mallas extremas alcanzaron 3.51\,\% en agua salada, 3.59\,\% en músculo y aproximadamente 0.26\,\% en concreto. La variación del orden aparente y del modelo de ajuste limitó la extrapolación al continuo.
\item El balance local a 5 px/mm produjo fracciones disipadas de 81.17\,\% en agua salada, 74.53\,\% en músculo y 30.38\,\% en concreto. En vacío la conductividad nula impuso pérdida de Joule cero. Los residuos finales quedaron por debajo de 0.055\,\% de la entrada total. La interpretación se refirió a los anillos seleccionados, con mayor variación temporal en vacío.
\item La fracción disipada de agua salada presentó una sensibilidad de 0.4181\,\% respecto de la malla de 10 px/mm. Las potencias absolutas cambiaron cerca del 30\,\%, por lo que su estabilidad excedió el umbral interno del 5\,\%. La resolución común resultó defendible para las tendencias y magnitudes relativas documentadas, con las limitaciones específicas de cada medio.
\item La pendiente aparente permitió formular un término de pérdida radial relativa para las configuraciones estudiadas. Su transferencia a modelos predictivos de propagación quedó condicionada a una validación independiente y a la calibración de la fuente. La penetración a través de interfaces y la eficiencia global de la antena requirieron verificaciones adicionales.
\end{enumerate}
